\section{Discussion}
\label{sec:discussion}

\para{Attention geometry as a failure signal.}
Our results provide the first empirical evidence that geometric properties of attention circuits contain predictive signal about prompt-level performance.
The significant correlations---particularly the Frobenius norm ($\spearman = +0.474$) and cross-layer variance ($\spearman = -0.472$)---support the hypothesis that geometric instability in attention patterns flags prompt failures.
These findings are consistent with prior work linking attention variance to Ricci curvature and model robustness \citep{park2024geometry}, and with the finding that mid-layer attention heads are both unstable and functionally important \citep{bali2026quantifying}.

\para{Why logit confidence fails.}
The complete failure of the logit-confidence baseline (0\% accuracy) is noteworthy.
Maximum softmax probability, a standard uncertainty measure \citep{hendrycks2017baseline}, cannot distinguish prompt formats that produce 71\% accuracy from those at chance level.
This suggests that the model is equally ``confident'' in its predictions regardless of whether the prompt format enables correct classification.
The failure of output-based heuristics motivates the use of internal geometric features, which capture structural properties of the computation rather than just the output distribution.

\para{Interpretability of geometric features.}
The correlations we observe have intuitive interpretations.
Higher Frobenius norm indicates stronger attention signals, suggesting the model is more effectively attending to task-relevant tokens.
Lower cross-layer variance indicates consistent attention behavior across the network, suggesting stable information flow.
Higher late-layer entropy suggests richer contextual integration in the final layers, where task-relevant representations are typically formed \citep{wang2022interpretability}.
These interpretations connect attention geometry to the mechanistic understanding of how prompts enable (or hinder) task performance.

\para{Practical implications.}
Geometric features require only a single forward pass to compute, making them a cheap complement to expensive evaluation-based prompt selection.
A practitioner could extract attention patterns for a small set of representative inputs, compute geometric features, and screen out formats with high cross-layer variance or low Frobenius norm before committing to full evaluation.
However, the current 60\% classification accuracy is insufficient for production use.

\subsection{Limitations}
\label{sec:limitations}

\para{Small sample size.}
The primary limitation is the small number of prompt formats ($n = 20$).
With only 20 data points, classifier generalization is poor and correlation confidence intervals are wide (several cross zero).
A production study using the \textsc{FormatSpread} grammar of \citet{sclar2023quantifying} to generate 100+ formats would provide much stronger statistical power.

\para{Single model.}
We evaluate only \gptsmall, the most studied model for mechanistic interpretability.
Its 12-layer, 12-head architecture may produce simpler attention geometry than larger models.
Whether the same geometric features are predictive in models like Llama-3 or GPT-4 remains an open question.

\para{Single task.}
\ssttwo is a relatively simple binary classification task.
Harder tasks such as multi-hop question answering or mathematical reasoning may exhibit different geometric signatures of failure.
Multi-task evaluation is needed to establish generality.

\para{Feature aggregation.}
Averaging geometric features across examples and heads discards fine-grained information.
Per-example prediction---predicting whether a specific (prompt, input) pair will fail---would be more powerful but requires a different experimental design with many more data points.

\para{No causal validation.}
Our correlations do not establish that geometry \emph{causes} failure.
Interventional experiments (\eg steering attention to modify geometric properties and measuring the effect on accuracy) would be needed to establish a causal link.

\para{Label word confound.}
Some accuracy variation may reflect label word frequency in \gptsmall's training data rather than prompt format per se.
Formats using ``positive/negative'' consistently outperform those using ``good/bad'' or ``great/terrible.''
Future work should disentangle label word effects from format effects.
